\section{The Newman boundary and exact positive heat matrices}

Use the normalization of the supplied workbench22,
\[
 H_t(z)=\int_0^\infty e^{tu^2}\Phi(u)\cos(zu)\,du,
 \qquad H_0(z)=\tfrac18\xi(1/2+iz/2),
 \qquad \partial_tH_t=-\partial_z^2H_t.
\]
The kernel $\Phi$ is positive and decays doubly exponentially. The
de Bruijn--Newman threshold property is imported as a known theorem:
all zeros of $H_t$ are real exactly when $t\ge\Lambda$. Rodgers--Tao,
arXiv:1801.05914v5, Theorem 1.1, proves $\Lambda\ge0$; A gives
$\Lambda\le0$. We use their intersection $\Lambda=0$, without claiming
that either theorem supplies simplicity at the endpoint.

\begin{theorem}[NN-A006: endpoint collision or escape to infinity]
Under A, every zero of $H_t$ is real and simple for $t>0$.
The zeros of $H_0$ are real, with their possible multiplicities retained.
If $t_j<0$, $t_j\uparrow0$, and $z_j$ is a nonreal zero of $H_{t_j}$,
then every finite accumulation point of $z_j$ is a multiple real zero
of $H_0$. In particular, if all zeta zeros are simple, such nonreal zeros
escape every compact subset of $\mathbb C$.

At a real multiplicity-$m$ zero $(t_0,z_0)$ the local branches are
\[
 z_j(t)=z_0+\sqrt{2(t-t_0)}\,h_{m,j}+O(|t-t_0|),
\]
where $h_{m,j}$ are the distinct real roots of the probabilists' Hermite
polynomial $\operatorname{He}_m$. For $m\ge2$, a real collision can occur only
at $t_0\le0$ under A. Replacing this by $t_0<0$ requires simplicity at
$t=0$, which A does not state.
\end{theorem}
\begin{proof}
The integral is jointly entire in $(t,z)$ locally, by the kernel decay,
so all compact-set convergences and derivatives used here are uniform.
Write $H_{t_0}(z)=a_m(z-z_0)^m+O((z-z_0)^{m+1})$.
The PDE gives
$\partial_t^\ell H=(-1)^\ell\partial_z^{2\ell}H$.
Set $t-t_0=u^2/2$ and $z-z_0=uy$; divide by $a_mu^m$.
The resulting analytic function equals at $u=0$
\[
 \sum_{\ell=0}^{\lfloor m/2\rfloor}
 \frac{m!(-1)^\ell y^{m-2\ell}}
 {2^\ell\ell!(m-2\ell)!}=\operatorname{He}_m(y).
\]
Gaussian differentiation and Rolle's theorem give $m$ real roots;
the Hermite differential equation excludes repeated roots. The implicit
function theorem at each root, with Rouche's theorem for completeness,
gives the stated expansion. For negative real $t-t_0$, at least two
of the nonzero Hermite roots give nonreal branches. If $t_0>0$, nearby
negative increments still have $t>0$, contradicting $\Lambda=0$.
This proves simplicity at all positive times.

If $z_j$ has a finite accumulation point $z_*$, uniform convergence
implies $H_0(z_*)=0$, and A makes $z_*$ real. If that zero were simple,
the implicit function theorem would give a unique local zero for all
small real $t$. Reality of the coefficients means its conjugate is
another local zero, so uniqueness forces the branch to be real. This
contradicts the nonreal $z_j$. Compactness gives the final escape claim.
\end{proof}

The theorem gives an exact dichotomy and no height law. In particular,
$\Lambda=0$ alone does not prove an exponential collision front or an
independent-pair collision approximation. In the completed arithmetic
coordinate $s=1/2+iz/2$, the double-zero law becomes
$s-s_0=\pm i\sqrt{(t-t_0)/2}+O(|t-t_0|)$; the normalization cannot
be omitted when comparing gap-based heuristics. The geometric local
model is a ramified double cover of the parameter disc (higher
Hermite branching for $m>2$), not a family with unramified individual
zero sections at the collision.

\begin{theorem}[NN-A009: strict heat positivity and determined moments]
Let the complete arithmetic expression $Z(a)$ be the one in
workbench22, equation \texttt{eq:rh-connes-arithmetic-heat}, including
both pole terms, the archimedean integral, and the prime-power sum.
Under A it has the exact representation
\[
 Z(a)=\sum_\rho e^{-a\gamma_\rho^2},\qquad a>0,
\]
with nontrivial zeros counted with multiplicity. Consequently:
\begin{enumerate}
\item $(-1)^kZ^{(k)}(a)>0$ for every $k\ge0$, and
$(\log Z)''(a)>0$.
\item If $a_1,\ldots,a_d>0$ are distinct, the matrix
$[Z(a_i+a_j)]_{i,j=1}^d$ is positive definite.
\item For every $d\ge1$, the matrix
$[(-1)^{i+j}Z^{(i+j)}(a)]_{i,j=0}^{d-1}$ is positive definite.
\item The positive finite measure
$\mu_a=\sum_\rho e^{-a\gamma_\rho^2}\delta_{\gamma_\rho^2}$
is uniquely determined among positive measures on $[0,\infty)$ by its
moments $(-1)^kZ^{(k)}(a)$.
\end{enumerate}
No simplicity assumption is needed.
\end{theorem}
\begin{proof}
The supplied unrestricted explicit formula is
$Z(a)=\sum_\rho e^{a(\rho-1/2)^2}$, absolutely and locally uniformly
convergent with all derivatives. A replaces $\rho-1/2$ by $i\gamma$.
There are infinitely many distinct positive values $\gamma^2$, since
an entire function has finitely many zeros in each compact set and
$\xi$ has infinitely many zeros. None is zero: $\zeta(1/2)\ne0$,
for example from the alternating eta series.
Termwise differentiation proves strict complete monotonicity.
After dividing the weights by $Z(a)$,
$Z''/Z-(Z'/Z)^2$ is the variance of $\gamma^2$, and is strictly
positive because its support has more than one point.

For the first matrix the quadratic form is
$\sum_\rho|\sum_i c_i e^{-a_i\gamma^2}|^2$.
A nonzero real linear combination of $d$ exponentials with distinct
real exponents has at most $d-1$ distinct real zeros: divide by one
exponential, differentiate, and use Rolle's theorem inductively on
any chosen finite interval containing the zeros. If a complex linear
combination vanished at infinitely many points, its real and imaginary
parts would each vanish identically by this fact; independence of the
exponentials then makes every coefficient zero. Thus the quadratic form
vanishes only at $c=0$. The second matrix has quadratic form
$\sum_\rho e^{-a\gamma^2}|\sum_{i=0}^{d-1}c_i\gamma^{2i}|^2$.
A nonzero polynomial has finitely many roots, proving strictness.

For determinacy, $\int e^{(a/2)x}\,d\mu_a=Z(a/2)<\infty$.
Any positive measure $\nu$ with the same moments has the same finite
exponential moment by monotone convergence of its nonnegative Taylor
series. Its moment generating function therefore agrees with that of
$\mu_a$ in $|z|<a/2$, by dominated termwise integration. Both Laplace
transforms are holomorphic on $\Re z<a/2$; the identity theorem gives
equality on that half-plane, hence equality of characteristic functions
on the imaginary axis. Fourier uniqueness for finite measures yields
$\nu=\mu_a$.
\end{proof}

The finite matrices are genuine Gram matrices of the displayed functions
in $L^2(\mu)$ or $L^2(\mu_a)$. An approximate matrix
$\widetilde G$ with a certified operator-norm error $\epsilon$ obeys
$\lambda_{\min}(\widetilde G)\ge\lambda_{\min}(G)-\epsilon$.
Under A any rigorously certified negative exact value is impossible;
an observed negative value must therefore fail either its error
enclosure, its admissibility, or the retained complete formula.
Strictness gives no uniform positive lower bound as the matrix size,
parameters, or conditioning vary. This is an analytic Hilbert-space
Gram realization. It is not a finite-dimensional pure Weil object,
and the tautological diagonal operator with spectrum $\gamma^2$
does not construct a new arithmetic Hilbert--Polya operator.
